% AUTOMATISCH ERZEUGT aus dem SysML-Modell mit tools/generate_docs.py.
% Nicht von Hand aendern: Inhalte im Modell pflegen, Layout im Kapitel.
TC-MAP-001 & Grundkarte laden & REQ-MAP-001 & Karte ist verfuegbar und wird von Lokalisierung und Navigation referenziert. \\
TC-MAP-002 & Karteninformation trennen & REQ-MAP-003 & Temporaere Hindernisse werden lokal beruecksichtigt, ohne die Langzeitkarte automatisch zu veraendern. \\
TC-LOC-001 & Lokalisierung in Karte & REQ-LOC-001 & Roboterpose wird innerhalb einer noch festzulegenden Toleranz geschaetzt. \\
TC-NAV-001 & Autonomes Zielanfahren & REQ-NAV-001 & Roboter erreicht das Ziel ohne manuelle Fahrzeugsteuerung. \\
TC-NAV-002 & Raeume-Flur-Tuer-Test & REQ-NAV-004, REQ-NAV-005 & Roboter durchfaehrt mindestens zwei Raeume, einen Flurabschnitt und eine offene Tuer autonom. \\
TC-OBS-001 & Hindernis frontal & REQ-MAP-002, REQ-NAV-002, REQ-SAFE-001 & Hindernis wird erkannt; Roboter stoppt oder faehrt sicher ausweichend. \\
TC-OBS-002 & Blockierter Fahrweg & REQ-NAV-003, REQ-SAFE-001, REQ-STA-001 & Roboter stoppt und meldet Blockade. \\
TC-UI-001 & Web-App-Auftrag & REQ-UI-001, REQ-MIS-001 & Start und Ziel werden ueber die Web-App gewaehlt; der Roboter nimmt den Auftrag an oder lehnt ihn mit nachvollziehbarem Status ab. \\
TC-UI-002 & Grundbefehle & REQ-UI-002 & Status, betrieblicher Stopp und Rueckkehr zur Ladestation erzeugen nachvollziehbare Systemreaktionen. \\
TC-PWR-002 & Unterrichtsprofil 2 h & REQ-PWR-002 & Roboter erreicht ein definiertes Demonstrationsprofil von ca. 2 h ohne Zwischenladen. \\
TC-PERF-001 & Nominalgeschwindigkeit & REQ-PERF-001 & Fahrprofil wird mit ca. $0{,}5\,\mathrm{m/s}$ innerhalb definierter Toleranz ausgefuehrt. \\
TC-PWR-001 & Niedriger Energiezustand & REQ-PWR-001 & Warnung bei niedrigem Ladezustand wird erzeugt und gemeldet. \\
TC-STA-001 & Missionsstatus melden & REQ-STA-001, REQ-MIS-002 & Statuswechsel werden korrekt ausgegeben. \\
TC-TRC-001 & Traceability-Review & REQ-ARCH-001, REQ-ARCH-002 & Jede MVP-Anforderung ist mit Funktion, Element, Schnittstelle und Test verknuepft; MVP-Elemente verweisen auf Zielsystem-Elemente. \\
TC-UI-004 & Abbruch und Quittierung & REQ-UI-002, REQ-MIS-002 & Abbruch fuehrt nach gestoppt, der Auftrag ist verworfen; erst die Quittierung fuehrt nach bereit, aus Fehler und sicherem Zustand nur bei erfuellten Voraussetzungen (DEC-015). \\
TC-DIAG-001 & Diagnosemeldung & REQ-DIAG-001 & Bei simuliertem Fehler enthaelt die Meldung Ursache, betroffene Funktion und Zeitstempel. \\
TC-MNT-001 & Wartungsreview & REQ-MNT-001 & Zugang zu zentralen Komponenten ohne unverhaeltnismaessige Demontage; erst mit mechanischem Entwurf ausfuehrbar. \\
TC-TRN-001 & Kistentransport & REQ-TRN-001 & Manuell aufgesetzter Werkstuecktraeger erreicht das Ziel; Nutzlastdaten und Zieltoleranz vor Ausfuehrung festlegen. \\
TC-TRN-002 & Lage der Kiste & REQ-TRN-002 & Werkstuecktraeger bleibt im vereinbarten Fahrprofil innerhalb der zulaessigen Lage; Grenzwerte offen. \\
TC-SAFE-002 & Person im Fahrweg & REQ-SAFE-002 & Reaktion auf eine Person im Fahrweg; Bewertungskriterium erst mit Safety-Konzept. \\
TC-UI-003 & Pause und manuelle Fortsetzung & REQ-UI-002, REQ-MIS-002 & Betrieblicher Stopp pausiert den Auftrag, Bewegung ist gestoppt, Auftrag bleibt erhalten; Fortsetzung nur bei manueller Freigabe und erfuellten Voraussetzungen (DEC-013). \\
